% !TeX root = ../main.tex
\section*{Conclusión}
La primera resurrección no requiere imaginar que los santos recuperan sus cuerpos mil años antes del juicio universal.
Las escenas de la guerra final recapitulan un mismo desenlace; el contexto inmediato presenta almas de testigos que han muerto; y las promesas anteriores del Apocalipsis asocian su muerte fiel con vida, reposo, victoria y preservación frente a la segunda muerte.
A la luz de esas conexiones, el lenguaje de resurrección revela el significado celestial de lo que en la tierra parecía una derrota.

Las preguntas que articulan la visión reciben así una respuesta conjunta.
¿Dónde reinan? En la esfera celestial anticipada por los tronos de Daniel y por el altar bajo el cual Juan había contemplado a los mártires.
¿Quiénes reinan? Los que perseveran en la fidelidad al Cordero, figurados de manera culminante por quienes derramaron su sangre.
¿Qué significa que \enquote{vivieron}? Que la muerte corporal no interrumpió su comunión con Cristo, sino que los introdujo en la manifestación gloriosa de su reino.
¿Y por qué \enquote{primera}? Porque esta bendición pertenece todavía al orden anterior a la consumación y anticipa, sin reemplazarla, la resurrección final.

La regeneración es el comienzo indispensable de esa vida en Cristo; el estado intermedio, su realización victoriosa después de la muerte; y la resurrección corporal, su perfección definitiva.
Por eso, en la paradoja de Juan, la muerte de los santos llega a ser resurrección, mientras que la resurrección de los impíos desemboca en la segunda muerte.
No son imágenes vacías: son la esperanza cristiana contemplada desde el trono del Cordero.
